\documentclass[11pt,a4paper]{article}
\usepackage[margin=25mm]{geometry}
\usepackage[T1]{fontenc}
\usepackage{lmodern,amsmath,amssymb,booktabs,longtable,array,hyperref}
\hypersetup{colorlinks=true,linkcolor=blue,urlcolor=blue}
\setlength{\emergencystretch}{3em}
\setlength{\parindent}{0pt}\setlength{\parskip}{6pt}
\title{Implementation Report v2\\Real-Data Reward Alignment for Portfolio Allocation}
\author{Matin Tavakoli}
\date{27 September 2026}
\begin{document}\maketitle
\begin{abstract}
This report documents the historical-data extension of the synthetic reward-design study. It explains the data pipeline, portfolio objectives, optimisation routines, experimental controls, recorded results and verification. The historical implementation tests whether alternative rewards select the intended fitted mean--variance allocation on the same estimated state. It also measures subsequent investment performance separately. Joint risk--cost calibration and an adapted mean--variance policy iteration method recover the fitted reference to numerical precision; the historical evidence does not establish superior realised performance. The synthetic implementation remains the v1 stage. This v2 report describes the additional real-data modules and their precise scope.
\end{abstract}
\setcounter{tocdepth}{1}
\begingroup\setlength{\parskip}{0pt}\tableofcontents\endgroup
\newpage
\section{What changed from the synthetic implementation}
Version 1 used specified return models and known analytical policies to isolate reward-design error. Its controlled setting made it possible to distinguish proved identities from numerical learning approximations. Version 2 replaces simulated observations with archived monthly US industry returns, adds lagged estimation of expected returns and covariance, imposes long-only portfolio constraints, and accounts for proportional turnover costs.

The real-data target is a fitted one-period allocation, not a known optimal policy for the true historical market. A rolling estimate does not identify the full conditional law of returns or a terminal-wealth equilibrium under dependence. The extension therefore tests implementation fidelity conditional on a specified model, then asks whether that fidelity helps realised performance. Those questions are deliberately recorded separately.

The principal new modules are \path{mvrl/historical.py} and \path{mvrl/historical_learning.py}. The first handles data, exact reward comparisons, portfolio paths, uncertainty and robustness. The second measures finite-budget numerical learning error. No synthetic strategy code was replaced by this extension.

\section{Data and the information boundary}
\subsection{Source and sample}
The source is Kenneth French's Data Library: ten value-weighted industry portfolio monthly total returns and the Fama--French risk-free series. The original files are \path{10_Industry_Portfolios_CSV.zip} and \path{F-F_Research_Data_Factors_CSV.zip}, retained under \path{data/raw/}. They were retrieved on 26 September 2026. The archived vintage is August 2026; SHA-256 hashes are recorded in \path{results/september26_downloads.json}.

The analysis retains January 1970 through December 2025, or 672 monthly observations. Portfolio evaluation starts January 1990 and ends December 2025: 432 monthly decisions. Earlier observations support estimation, including the longest 240-month window. The ten industry portfolios are research portfolios, not directly investable funds. Their internal constituent-rebalancing costs are not included.

The parser selects the value-weighted monthly block, recognises six-digit year--month rows, converts percentages to decimal returns, and rejects French missing-value sentinels. Asset and risk-free observations are joined by month. Missing joined data and duplicate dates are rejected. The data are a current revised vintage; lagging observations prevents use of future monthly returns but does not recreate point-in-time historical releases or constituent membership.

\subsection{Lagged estimates}
At month $t$, the primary window contains the preceding 120 months only. With excess asset returns $X_s=R_s-r_{f,s}\mathbf1$, the code calculates the arithmetic sample mean $m_t$ and the unbiased sample covariance $S_t$. The covariance used in every reward is
\[
\Sigma_t=0.9S_t+0.1\operatorname{diag}(S_t).
\]
There is no primary mean shrinkage. Positive definiteness is checked. All methods share these estimates and the same decision dates. The records store the first and last fitting month, so the lag boundary can be independently reconstructed.

\section{Financial objective and reward definitions}
\subsection{Reference allocation}
The feasible set is the ten-asset simplex $C=\{u:u_i\ge0,\ \sum_i u_i=1\}$. There is no cash decision or leverage. With pre-trade weights $p$, the target is
\[
F(u)=m^\top u-\varphi u^\top\Sigma u-\lambda\|u-p\|_1.
\]
The primary coefficients are $\varphi=5$ and $\lambda=0.001$. The latter means ten basis points per unit of L1 turnover. A transfer of weight $\delta$ from one industry to another contributes $2\delta$ to this measure. These coefficients apply to monthly decimal returns; changing frequency or units without recalibration changes their interpretation.

\subsection{Implemented comparator family}
Let $M=\Sigma+mm^\top$, $\nu=m^\top M^{-1}m$, $u^*=\arg\max_C F$, $y^*=m^\top u^*$ and $A=1+2\varphi y^*$. The code rejects a nonpositive $A$ for joint calibration.
\begin{longtable}{p{.24\textwidth}p{.68\textwidth}}
\toprule Method & Implemented definition \\\midrule
Target & $m^\top u-\varphi u^\top\Sigma u-\lambda\|u-p\|_1$.\\
Raw second moment & $m^\top u-\varphi u^\top M u-\lambda\|u-p\|_1$.\\
Unconstrained scalar & Replaces the raw quadratic coefficient with $\varphi(1-\nu)$; leaves the cost coefficient unchanged.\\
Risk-only & Uses $\varphi/A$ for the second moment but retains cost $\lambda$.\\
Joint calibration & Uses $\varphi/A$ for the second moment and $\lambda/A$ for turnover.\\
MVPI adaptation & Alternates a mean auxiliary with a constrained portfolio solve; does not require inserting the target allocation first.\\
Log minus variance & Window-average $\log(1+u^\top R_s)$, minus $\varphi u^\top\Sigma u$ and the linear turnover cost.\\
Log utility & Same logarithmic objective with zero variance coefficient.\\
Equal weight & Fixed one-tenth allocation in each industry, rebalanced monthly.\\
Minimum variance & Target objective with its expected-return term removed; risk scaling and turnover charge retained.\\\bottomrule
\end{longtable}

\subsection{Why risk and costs must be calibrated together}
For the joint reward $Q$, direct expansion gives
\[
A Q(u)=F(u)-\varphi(m^\top u-y^*)^2+\varphi(y^*)^2.
\]
At $u^*$, the additional negative square is zero. At every other feasible allocation it is nonpositive, while $F(u)\le F(u^*)$. Thus the reference also maximises $Q$. Strict concavity ensures uniqueness under the implemented covariance assumptions. This is a proved consequence of the established mean--variance auxiliary identity, not a newly invented general-purpose optimisation algorithm.

If only risk is rescaled, the right side gains $-2\varphi y^*\lambda\|u-p\|_1$. For positive reference expected return, that term increases the effective cost preference. It need not imply lower realised turnover over a whole path. The joint formula is a reference-dependent alignment certificate. MVPI provides a constructive comparison: starting at $y=0$, maximise $(1+2\varphi y)m^\top u-\varphi u^\top Mu-\lambda\|u-p\|_1$, then set $y=m^\top u$. The code uses tolerance $10^{-10}$ and a maximum of 200 iterations.

\section{Optimisation and numerical checks}
The quadratic solver introduces ten auxiliary variables satisfying $z_i\ge u_i-p_i$ and $z_i\ge p_i-u_i$, along with the simplex constraints. It minimises the negative reward using SciPy SLSQP. The initial scaling is 1,000, the stopping tolerance is $10^{-11}$ and the initial iteration limit is 300. Bounds on weights and auxiliary variables are zero to one.

A separate linear programme over the same epigraph polytope computes a first-order objective-gap certificate in unscaled units. For convex $f$ and feasible $x$, $C(x)=\max_{v\in C}\nabla f(x)^\top(x-v)$ bounds $f(x)-\min f$. In this expression $C$ denotes the extended feasible polytope, including the auxiliary variables. The certificate is independent of the SLSQP success flag.

If the initial solve reports failure or has a certificate above $2\times10^{-7}$, the solver retries at scales 10, 1 and 100, with tolerance $10^{-13}$ and up to 1,000 iterations. Nonfinite or infeasible retry candidates are rejected. A remaining excessive certificate raises an error. The largest recorded common-state certificate in the primary run is $4.2675\times10^{-9}$.

Logarithmic comparisons use a separate SLSQP objective with an analytical gradient, scale 1,000, tolerance $10^{-10}$ and up to 300 iterations. The quadratic certificate does not certify those logarithmic solutions. The report therefore distinguishes the stronger quadratic verification from the log solver's convergence check.

\section{Common-state auditing and realised paths}
At every evaluation month, the target strategy's pre-trade weights define the common audit state. All ten methods are solved on it. The saved Euclidean weight distance is $\|u^r-u^*\|_2$; fitted regret is $F(u^*)-F(u^r)$. These isolate the consequence of changing the reward on identical information and holdings.

Separately, each strategy solves its objective on its own pre-trade weights. If those weights differ from the common state, its realised allocation is recomputed. Its turnover is $T_t=\|u_t-p_t\|_1$, gross return is $u_t^\top R_t$, and net return is $u_t^\top R_t-\lambda T_t$. Relative holdings then drift as
\[
p_{t+1,i}=\frac{u_{t,i}(1+R_{t,i})}{1+u_t^\top R_t}.
\]
All strategies start equal-weighted. This is first-order proportional-cost accounting, not an exact self-financing execution or market-impact model. Wealth compounds net returns. The records retain both realised weights and common-state audit weights; these must not be interchanged in subsequent analysis.

\section{Results and what they establish}
The following table is generated directly from the saved primary summary. Fitted gap and realised MV are monthly basis points; allocation distance is Euclidean and Sharpe is annualised from monthly excess returns.
\begin{center}\small
\begin{tabular}{lrrrr}\toprule
Method & Distance & Fitted gap & Sharpe & Realised MV\\\midrule
Target & 0.00000 & 0.0000 & 0.570 & -8.497 \\
Raw second moment & 0.00843 & 0.0148 & 0.573 & -7.182 \\
Unconstrained scalar & 0.00540 & 0.0043 & 0.568 & -8.882 \\
Risk-only & 0.00510 & 0.0033 & 0.571 & -8.365 \\
Joint calibration & 0.00000 & 0.0000 & 0.570 & -8.497 \\
Equal weight & 0.43495 & 46.2626 & 0.622 & -11.893 \\
Minimum variance & 0.23843 & 9.0252 & 0.645 & 5.274 \\
MVPI adaptation & 0.00000 & 0.0000 & 0.570 & -8.497 \\
Log utility & 0.54848 & 39.7855 & 0.344 & -133.295 \\
Log minus variance & 0.01286 & 0.0287 & 0.572 & -7.123 \\
\bottomrule\end{tabular}\end{center}
Joint calibration and MVPI reproduce the fitted target to numerical precision. Raw, scalar and risk-only choices retain discrepancies. Logarithmic objectives represent different preferences, so disagreement with an MV target does not by itself invalidate them. The minimum-variance control has both better realised MV and higher Sharpe than the fitted target in the primary sample. Equal weight has higher Sharpe but lower realised MV than the target. These observations show why alignment, estimation quality and performance rankings must be separated.

Paired circular-block bootstrap intervals use 12-month blocks, 1,000 resamples and seed 8128. The same block indices are used for both return paths. For joint calibration minus the named comparator, the monthly MV differences and 95\% intervals, in basis points, are:
\begin{center}\begin{tabular}{lrr}\toprule Comparator & Difference & 95\% interval\\\midrule
Raw second moment & -1.315 & [-3.318, 0.343] \\
Unconstrained scalar & 0.385 & [-0.570, 1.351] \\
Risk-only & -0.131 & [-0.889, 0.634] \\
\bottomrule\end{tabular}\end{center}
The principal intervals contain zero. The bootstrap conditions on the saved strategy returns; it does not refit estimates within every resample or include hyperparameter-search uncertainty. The result is not a demonstration of realised investment superiority.

\section{Robustness and chronological checks}
Two grids produce 17 unique specifications: risk coefficients $\{2,5,10\}$ crossed with costs $\{0,0.001,0.0025\}$ at the primary window and shrinkage; and windows $\{60,120,240\}$ crossed with shrinkage $\{0,0.1,0.5\}$ at the primary risk and cost. The primary setting appears in both grids and is counted once. This is not the full Cartesian product of all four parameters.

Each robustness run includes target, raw, scalar, risk-only, joint, equal weight and minimum variance. MVPI and the two logarithmic comparisons are primary-run comparators; they are not rerun throughout the robustness grid. The maximum setting-mean joint distance is $1.5354\times10^{-7}$. Raw mean distances range from about 0.00157 to 0.03682. The 17 settings reuse the same observations and are not independent replications.

Saved chronological summaries cover 1990--2007, 2008--2015 and 2016--2025. They summarise the corresponding segments of the original paths, rather than retraining a new model or resetting portfolios at each boundary. A second asset universe was suggested in the protocol but is not part of the completed evidence.

\section{Finite-budget learning diagnostic}
The diagnostic reconstructs the 432 lagged common states from the primary experiment. Each state has its own ten softmax logits; there is no shared predictive network. PyTorch uses float64, one computation thread, Adam with learning rate 0.05, and initial logits drawn from $N(0,0.1^2)$. Seeds 0--4 are matched across target, raw and joint rewards. Checkpoints are 50, 250, 1,000 and 5,000 gradient steps. The minimised loss is minus 1,000 times the mean reward over the states.

No subsequent realised return enters a state's optimisation inputs. Nevertheless, this is an optimisation diagnostic, not evidence of held-out neural generalisation: each state has independent trainable parameters. Softmax cannot attain exact zero weights at finite logits, which can make boundary optima slower to approximate.

Let $u_r$ be the exact reward optimum and $\hat u_r$ the finite-budget output. The saved target decomposition is
\[
F(u^*)-F(\hat u_r)=\underbrace{F(u^*)-F(u_r)}_{\text{design gap}}+
\underbrace{F(u_r)-F(\hat u_r)}_{\text{signed transfer}}.
\]
The second term can be negative. Own-reward suboptimality $Q(u_r)-Q(\hat u_r)$ is recorded separately and must not replace it in the equality. The saved file contains 60 aggregate records: three methods times five seeds times four budgets. At 5,000 steps, mean own-reward gaps are approximately $4\times10^{-6}$, larger than the primary raw design gap $1.4801\times10^{-6}$. This preliminary evidence separates learning effort from the discrepancy remaining at exact optimisation. It is not a PPO, DDPG, TD3 or CNN/WaveNet reproduction.

\section{Published comparisons and attribution}
Two publications supply implemented adaptations. Zhang, Liu and Whiteson's \emph{Mean-Variance Policy Iteration for Risk-Averse Reinforcement Learning} (AAAI, 2021; DOI \url{https://doi.org/10.1609/aaai.v35i12.17302}) supplies the auxiliary method. Jiang, Olmo and Atwi's \emph{High-dimensional multi-period portfolio allocation using deep reinforcement learning} (International Review of Economics and Finance 98, 2025, 103996; DOI \url{https://doi.org/10.1016/j.iref.2025.103996}), Eq. (8), supplies the log-minus-variance preference. Its Eq. (19) includes restrictions which are replaced here by the common hard simplex; its cost treatment is adapted to the shared linear convention.

Recent literature comparisons also include Jiang et al. (Global Finance Journal, 2024), Gollart and Okhrin (Annals of Operations Research, online 2025), Simon, Weibels and Zimmermann (Management Science, 2026), and Cong, Tang and Wang (NBER Working Paper 35195, May 2026). Their full learning systems are not empirical baselines in this implementation. Exact bibliographic records, locators and access notes are retained in the project literature ledger. The strongest theoretical increment remains the synthetic fixed-horizon mismatch coefficient; its priority is not certified by the historical implementation.

\section{Verification, outputs and reproduction}
\subsection{Tests and independent reconstruction}
The saved research verification comprises 31 passing pytest tests and ten original checks. Four historical regression tests cover constrained joint/MVPI agreement, the risk-only cost-distortion identity, an independent dense two-asset grid comparison with the quadratic certificate, and parsing/lagged moment behaviour. The independent historical checker reconstructs 4,320 primary strategy decisions, checks file hashes, lag dates, feasibility, turnover, net-return accounting, drift, common-state regret and summary consistency. It also checks all 17 robustness summaries and the 60-record signed learning decomposition.

This checker reconstructs accounting from records rather than rerunning every solver, and it imports the common data/moment helpers. It should not be described as an entirely independent implementation of the whole pipeline. The two-asset grid test supplies a separate numerical check of a small optimisation problem.

\subsection{Saved artefacts}
\begin{longtable}{p{.44\textwidth}p{.49\textwidth}}\toprule Path & Purpose\\\midrule
\path{results/historical/primary.json} & Per-method monthly paths and common-state audit records.\\
\path{results/historical/summary.json} & Primary metrics, paired intervals and chronological summaries.\\
\path{results/historical/robustness.json} & The 17 setting summaries; separate setting JSON files retain raw paths.\\
\path{results/historical/learning.json} & Sixty finite-budget diagnostic records.\\
\path{scripts/check_historical.py} & Reconstruction checks on the archived results.\\
\path{scripts/make_historical_figures.py} & Reproducible figures from saved results.\\
\path{notes/REAL_DATA_PROTOCOL.md} & Recorded experimental design and scope.\\
\path{notes/CONSTRAINED_CALIBRATION.md} & Identity, interpretation and prior-art attribution.\\\bottomrule
\end{longtable}

\subsection{Run order}
From the canonical project directory, install the recorded dependencies and set \path{PYTHONPATH=.}. PyTorch is required for the learning diagnostic. The saved core and learning environment files record the versions used. Execute:
\begin{verbatim}
python -m pytest -q
python -m mvrl.verify
python -m mvrl.historical --robust
python -m mvrl.historical_learning
python scripts/make_historical_figures.py
python scripts/check_historical.py
\end{verbatim}
The historical runner reuses saved primary and setting files when they exist. For a genuinely fresh experiment after changing data or objectives, first move the historical-results directory to a dated archive; do not combine new code with old cached results and describe that as a new run. The raw ZIPs, hashes, parameter settings, seeds and environment records define the reported study.

\newpage
\section{Result status and scope}
\begin{longtable}{p{.48\textwidth}p{.45\textwidth}}\toprule Result & Status\\\midrule
Joint constrained risk--cost identity and risk-only distortion & Proved, as consequences of the established auxiliary representation.\\
Historical common-state alignment across 17 settings & Numerically verified for the archived data and settings.\\
Finite-budget softmax learning comparisons & Preliminary; five seeds, no predictive generalisation claim.\\
Higher realised performance from calibration & Not established; principal paired intervals include zero.\\
True historical dependent terminal equilibrium & Open; not supplied by rolling one-period estimation.\\
Full published RL architecture replication & Not implemented in this extension.\\
Point-in-time investable-universe and execution validation & Open extension.\\\bottomrule
\end{longtable}
The completed implementation makes one practical capability available: given a fitted preference and a proposed reward, it can measure their policy disagreement on real historical states, calibrate the risk--cost reward when the identity applies, and separately report optimisation error and investment outcomes. It does not turn fitted preference alignment into a guarantee of superior trading performance.

\end{document}
